\documentclass[11pt,reqno]{amsart}

\usepackage{amsmath,amssymb,amsthm,mathtools}
\usepackage{hyperref}
\usepackage{microtype}

\newtheorem{theorem}{Theorem}[section]
\newtheorem{proposition}[theorem]{Proposition}
\newtheorem{corollary}[theorem]{Corollary}
\newtheorem{lemma}[theorem]{Lemma}
\theoremstyle{definition}
\newtheorem{definition}[theorem]{Definition}
\theoremstyle{remark}
\newtheorem{remark}[theorem]{Remark}

\DeclareMathOperator{\supp}{supp}
\DeclarePairedDelimiter{\abs}{\lvert}{\rvert}
\DeclarePairedDelimiter{\norm}{\lVert}{\rVert}

\title{The Half-Kernel of the Zeta Spectral Process}
\date{\today}

\begin{document}
\maketitle

\section{Setup}

Let \(\theta\) be the Riemann--Siegel theta function, \(Z(t)=e^{i\theta(t)}\zeta(\tfrac{1}{2}+it)\)
the Hardy \(Z\)-function, and \(u=\theta(t)\) the phase variable. Define the
time-changed process
\[
  h(u) := \frac{Z(t(u))}{\sqrt{\theta'(t(u))}},
  \qquad t(u) = \theta^{-1}(u),
\]
which is weakly stationary on \([\theta(T_0),\infty)\) with \(T_0=200\).
Its Cram\'{e}r spectral representation is
\[
  h(u) = \int_{-1}^{1} e^{i\nu u}\,dM(\nu),
  \qquad \mathbb{E}\abs{dM(\nu)}^2 = S(\nu)\,d\nu + w\,\delta_{-1},
\]
where \(\supp S \subseteq [-1,1]\) and \(S(\nu)\,d\nu\) is the absolutely
continuous part of the spectral measure \(\nu\).

The spectral density \(S\) is given by the limit
\begin{equation}\label{eq:S}
  S(\nu) = \lim_{T\to\infty} \abs{K_T(\nu)}^2,
  \qquad
  K_T(\nu) = \frac{1}{2\pi}\int_{U_0}^{U_T} h(u)\,e^{-i\nu u}\,du,
\end{equation}
which in terms of the Riemann--Siegel expansion
\[
  h(u) = \sum_{n=1}^{N(t(u))}
    \frac{e^{i(u-t(u)\log n)}+e^{-i(u-t(u)\log n)}}{\sqrt{n}\,\sqrt{\theta'(t(u))}}
  + \text{lower order}
\]
takes the form
\begin{equation}\label{eq:S-double}
  S(\nu) = \lim_{T\to\infty} \frac{1}{4\pi^2}
  \sum_{m,n=1}^{N_T} \frac{1}{\sqrt{mn}}
  \iint \frac{e^{i\Psi_m(s)-i\Psi_n(r)}}
             {\sqrt{\theta'(t(s))\,\theta'(t(r))}}
  \,e^{-i\nu(s-r)}\,ds\,dr,
\end{equation}
where \(\Psi_n^\pm(u,\nu) = \pm(u-t(u)\log n) - \nu u\) are the
Riemann--Siegel phases. The off-diagonal \(m\neq n\) cross terms in
\eqref{eq:S-double} encode the arithmetic pair-correlation structure of
the integers through the irrationality of \(\log m/\log n\).

\section{The Half-Kernel}

\begin{definition}[Spectral factor / half-kernel]\label{def:halfkernel}
The \emph{half-kernel} \(k\in L^2(\mathbb{R})\) of the zeta spectral process
is the inverse Fourier transform of the amplitude \(\sqrt{S(\nu)}\):
\[
  k(u) := \frac{1}{2\pi}\int_{-1}^{1} \sqrt{S(\nu)}\,e^{i\nu u}\,d\nu.
\]
It satisfies the factorization identity
\[
  (k * \tilde{k})(u) = C_h(u),
  \qquad \tilde{k}(u) := \overline{k(-u)},
\]
where \(C_h(\tau) = \mathbb{E}[h(u+\tau)\overline{h(u)}]\) is the
autocovariance of \(h\), and the moving-average representation
\[
  h(u) = \int_{\mathbb{R}} k(u-v)\,dW(v)
\]
holds in \(L^2(\Omega,\mathbb{P})\), where \(dW\) is the white-noise measure
on \(\mathbb{R}\).
\end{definition}

\begin{theorem}[Half-kernel as phase-flattened convolution]\label{thm:main}
Let \(\phi(\nu) = \arg \hat{h}(\nu)\) be the spectral phase of \(h\), where
\(\hat{h}(\nu) = \lim_{T\to\infty} K_T(\nu)\). Define the
\emph{phase-flattening kernel}
\[
  \kappa(\xi) := \frac{1}{2\pi}\int_{-1}^{1} e^{i\nu\xi}\,e^{-i\phi(\nu)}\,d\nu.
\]
Then the half-kernel satisfies
\begin{equation}\label{eq:main}
  k(u) = (h * \kappa)(u) = \int_{\mathbb{R}} h(s)\,\kappa(u-s)\,ds,
\end{equation}
as an equality in \(L^2(\mathbb{R})\). Equivalently,
\begin{equation}\label{eq:main-RS}
  k(u) = \sum_{n=1}^{\infty} \frac{1}{\sqrt{n}}
  \int_{\mathbb{R}}
  \frac{e^{i(s-t(s)\log n)}+e^{-i(s-t(s)\log n)}}{\sqrt{\theta'(t(s))}}
  \,\kappa(u-s)\,ds,
\end{equation}
a superposition over all integers \(n\geq 1\) of phase-flattened
translates of the Riemann--Siegel amplitudes.
\end{theorem}

\begin{proof}
Write \(\sqrt{S(\nu)} = \abs{\hat{h}(\nu)} = \hat{h}(\nu)\,e^{-i\phi(\nu)}\).
Then
\[
  k(u) = \frac{1}{2\pi}\int_{-1}^{1} \hat{h}(\nu)\,e^{-i\phi(\nu)}\,e^{i\nu u}\,d\nu.
\]
Since \(\hat{h}\in L^2([-1,1])\) by the Cram\'{e}r isomorphism and
\(e^{-i\phi(\nu)}\) is bounded (of modulus \(1\)), the product
\(\hat{h}(\nu)e^{-i\phi(\nu)}e^{i\nu u}\in L^1([-1,1])\). Substituting
\[
  \hat{h}(\nu) = \int_{\mathbb{R}} h(s)\,e^{-i\nu s}\,ds
\]
(as an \(L^2\) Fourier transform) and exchanging the \(\nu\)-integral
with the \(s\)-integral by the \(L^2\) form of Fubini's theorem
(justified because \(\hat{h}\in L^2([-1,1])\) and the \(\nu\)-domain
is finite), we obtain
\[
  k(u)
  = \int_{\mathbb{R}} h(s)
    \left(\frac{1}{2\pi}\int_{-1}^{1} e^{i\nu(u-s)}\,e^{-i\phi(\nu)}\,d\nu\right)ds
  = (h * \kappa)(u).
\]
Equation \eqref{eq:main-RS} follows by substituting the Riemann--Siegel
expansion of \(h\) and using linearity of convolution.
\end{proof}

\begin{remark}
The kernel \(\kappa\) is real and even because \(h\) is real-valued,
which forces \(\phi(-\nu)=-\phi(\nu)\) (odd spectral phase), making
\(\kappa(\xi)=\frac{1}{\pi}\int_0^1\cos(\nu\xi-\phi(\nu))\,d\nu\) manifestly real.
Consequently \(k\) is real, consistent with the half-kernel of a
real-valued stationary process.
\end{remark}

\begin{remark}[Arithmetic content]
Each term in \eqref{eq:main-RS} contributes oscillations at the
Riemann--Siegel frequencies \(\nu_n^\pm = \pm 1 \mp \log n/\theta'(t)\),
which are dense in \([-1,1]\) as \(n\) and \(t\) vary. The
phase-flattening kernel \(\kappa\) acts coherently across all \(n\),
stripping the phase while preserving the modulus \(\sqrt{S(\nu)}\). The
resulting \(k\) therefore encodes the full arithmetic of the integer
pairs \((m,n)\) through the off-diagonal terms of \eqref{eq:S-double},
in a single band-limited \(L^2\) function supported spectrally on
\([-1,1]\).
\end{remark}

\begin{remark}[Prolate connection]
The half-kernel \(k\in PW_{[-1,1]}\) (the Paley--Wiener space of
functions band-limited to \([-1,1]\)). The integral operator with kernel
\(C_h(u-v) = (k*\tilde{k})(u-v)\) restricted to a finite interval
\([-T,T]\) is a finite-section Toeplitz operator whose eigenfunctions
are prolate spheroidal wave functions with bandwidth parameter \(c=T\).
The weight \(\sqrt{S(\nu)}\) on the spectral side specifies how the
prolate eigenfunctions are weighted relative to the uniform (Lebesgue)
measure on \([-1,1]\).
\end{remark}

\end{document}
